%%%%%%%%%%%%%%%%%%%%%%%%%%%%%%%%%%%%%%%%%%%%%%%%%%%%%%%%%%%%%%
% CHAPTER 3
%%%%%%%%%%%%%%%%%%%%%%%%%%%%%%%%%%%%%%%%%%%%%%%%%%%%%%%%%%%%%%

\chapter{Dataset Collection}

\section{Introduction}

The quality of an NLP system largely depends on the quality, diversity, and suitability of the dataset used for training and evaluation. Dataset selection is therefore a critical stage in the development of any Natural Language Processing application. The selected dataset should adequately represent the target domain, language, and NLP task while ensuring sufficient quantity and quality of textual data.

This chapter describes the criteria used for selecting the \textit{DocSaarthi} dataset, its source, its characteristics, the annotation methodology followed, licensing information, and the justification for its use in the proposed project. The dataset is used to train and evaluate the Document Classification and Key Information Extraction stages of the proposed pipeline.

%%%%%%%%%%%%%%%%%%%%%%%%%%%%%%%%%%%%%%%%%%%%%%%%%%%%%%%%%%%%%%

\section{Dataset Selection Criteria}

The dataset satisfies the following selection criteria.

\begin{itemize}
    \item \textbf{Relevance to the selected NLP task:} Every record is a Hindi legal or financial document text pre-labelled with a document type and an intent, directly matching the Document Classification and Key Information Extraction stages of the proposed pipeline.
    \item \textbf{Availability:} The dataset is self-created and locally available in CSV format, ensuring unrestricted access throughout the project.
    \item \textbf{Sufficient number of records:} The dataset contains 400 records, adequate for training and evaluating classical machine learning baselines and fine-tuning lightweight transformer classifiers within the scope of a micro project.
    \item \textbf{Appropriate language:} All text is in Hindi (Devanagari script), directly matching the target output language of the translation stage of the proposed system.
    \item \textbf{Availability of annotations:} Each record carries two labels -- \texttt{Document\_Type} (Legal / Financial) and \texttt{Intent} (8 sub-classes) -- required for the classification stage.
    \item \textbf{Good data quality:} The dataset contains no missing values and no duplicate records (verified in Section~\ref{sec:datasetcharacteristics}).
    \item \textbf{Balanced class distribution:} The two primary classes (Legal and Financial) are perfectly balanced at 200 records each.
    \item \textbf{Appropriate for academic use:} As a self-created dataset, no external licensing restriction applies.
\end{itemize}

%%%%%%%%%%%%%%%%%%%%%%%%%%%%%%%%%%%%%%%%%%%%%%%%%%%%%%%%%%%%%%

\section{Dataset Source}

The dataset used in this project, \texttt{DocSaarthi\_Dataset.csv}, is a \textbf{self-created dataset}. It was synthetically generated by populating a small number of Hindi legal and financial document templates (service agreements, court proceeding notices, loan documents, insurance policies, tax notices, payment notices, and confidentiality clauses) with randomly varied entities -- party names, organization names, dates, monetary amounts, and durations -- to produce 400 distinct, realistic-looking legal and financial text records.

This approach was adopted because publicly available, freely licensed, sentence-level Hindi legal and financial text corpora with document-type and intent labels are scarce, and it allowed full control over class balance across both the Document\_Type and Intent labels. The self-created nature of the dataset is disclosed transparently, consistent with the annotation and licensing discussion in Sections~\ref{sec:annotationstrategy} and~\ref{sec:datasetlicensing}.

\begin{itemize}
    \item \textbf{Dataset file:} \texttt{DocSaarthi\_Dataset.csv}
    \item \textbf{Access date:} As used for this micro project, Academic Year 2026--2027
\end{itemize}

%%%%%%%%%%%%%%%%%%%%%%%%%%%%%%%%%%%%%%%%%%%%%%%%%%%%%%%%%%%%%%

\section{Dataset Description}

Table~\ref{tab:datasetdescription} summarizes the selected dataset.

\renewcommand{\arraystretch}{1.5}

\begin{longtable}{|p{5cm}|p{10cm}|}
\caption{Dataset Description}
\label{tab:datasetdescription}\\

\hline
\textbf{Parameter} & \textbf{Description} \\
\hline
\endfirsthead

\hline
\textbf{Parameter} & \textbf{Description} \\
\hline
\endhead

Dataset Name &
DocSaarthi\_Dataset
\\ \hline

Source &
Self-created (synthetically generated from Hindi legal/financial document templates)
\\ \hline

URL &
Not applicable (locally created dataset)
\\ \hline

Domain &
Legal and Financial documents
\\ \hline

NLP Task &
Document Classification (Legal/Financial), Intent Classification, Key Information Extraction
\\ \hline

Language(s) &
Hindi (Devanagari script)
\\ \hline

File Format &
CSV
\\ \hline

Dataset Size &
400 records, 4 columns (\texttt{ID}, \texttt{Text}, \texttt{Document\_Type}, \texttt{Intent})
\\ \hline

Training Split &
280 records (70\%) -- proposed split, to be applied during implementation
\\ \hline

Validation Split &
60 records (15\%) -- proposed split, to be applied during implementation
\\ \hline

Testing Split &
60 records (15\%) -- proposed split, to be applied during implementation
\\ \hline

License &
Self-created for academic use; no third-party licensing restrictions
\\ \hline

Version &
v1.0
\\ \hline

Date Accessed &
Academic Year 2026--2027
\\ \hline

\end{longtable}

%%%%%%%%%%%%%%%%%%%%%%%%%%%%%%%%%%%%%%%%%%%%%%%%%%%%%%%%%%%%%%

\section{Dataset Characteristics}
\label{sec:datasetcharacteristics}

Table~\ref{tab:datasetcharacteristics} presents the important characteristics of the dataset, computed directly from \texttt{DocSaarthi\_Dataset.csv}.

\renewcommand{\arraystretch}{1.5}

\begin{longtable}{|p{6cm}|p{9cm}|}
\caption{Dataset Characteristics}
\label{tab:datasetcharacteristics}\\

\hline
\textbf{Characteristic} & \textbf{Observation} \\
\hline
\endfirsthead

\hline
\textbf{Characteristic} & \textbf{Observation} \\
\hline
\endhead

Number of Records &
400
\\ \hline

Number of Classes (Document\_Type) &
2 (Legal, Financial)
\\ \hline

Number of Classes (Intent) &
8 (Agreement, Court Proceeding, Confidentiality, Finance, Loan, Payment, Tax, Insurance)
\\ \hline

Class Labels (Document\_Type distribution) &
Legal: 200 (50.0\%); Financial: 200 (50.0\%)
\\ \hline

Class Labels (Intent distribution) &
Agreement: 70; Court Proceeding: 65; Confidentiality: 65; Finance: 40; Loan: 40; Payment: 40; Tax: 40; Insurance: 40
\\ \hline

Average Sentence/Record Length (words) &
46.9 words
\\ \hline

Maximum Record Length (words) &
71 words
\\ \hline

Minimum Record Length (words) &
30 words
\\ \hline

Average Record Length (characters) &
262.1 characters
\\ \hline

Vocabulary Size &
1,196 unique whitespace-delimited Hindi tokens
\\ \hline

Language Distribution &
100\% Hindi (Devanagari script)
\\ \hline

Missing Values &
0 (no null values in any column)
\\ \hline

Duplicate Records &
0 (no duplicate \texttt{Text} entries)
\\ \hline

Special Characters &
Devanagari punctuation (।), currency symbol (₹), numerals, forward slashes in dates
\\ \hline

Encoding Format &
UTF-8
\\ \hline

\end{longtable}

%%%%%%%%%%%%%%%%%%%%%%%%%%%%%%%%%%%%%%%%%%%%%%%%%%%%%%%%%%%%%%

\subsection{Sample Records}

Table~\ref{tab:samplerecords} shows representative records from the dataset, one for each of four Intent categories.

\begin{table}[H]
\centering
\caption{Sample Dataset Records}
\label{tab:samplerecords}

\renewcommand{\arraystretch}{1.4}

\begin{tabular}{|c|p{8cm}|c|}
\hline

Record ID &
Text &
Label
\\

\hline

1 &
यह अनुबंध यह सेवा अनुबंध दिनांक 23/05/2025 को आर्यन ट्रेडर्स तथा नेहा शर्मा के बीच संपन्न हुआ। अनुबंध की अवधि 29 माह निर्धारित की गई है। कर्मचारी को ₹31717 मासिक पारिश्रमिक दिया जाएगा। ...
&
Legal / Agreement
\\

\hline

71 &
माननीय न्यायालय ने वाद संख्या 1757/2026 में कविता पाटिल बनाम कविता गुप्ता का मामला जिला न्यायालय में विचाराधीन है। अगली सुनवाई 08/09/2025 को निर्धारित की गई है। ...
&
Legal / Court Proceeding
\\

\hline

202 &
ऋण दस्तावेज़ के अनुसार भारतीय स्टेट बैंक ने अजय पाटिल को ₹564670 का ऋण 14/12/2025 को स्वीकृत किया। वार्षिक ब्याज दर 9.83\% है। मासिक ईएमआई ₹10509 निर्धारित की गई है। ...
&
Financial / Loan
\\

\hline

205 &
कविता वर्मा ने एचडीएफसी लाइफ की बीमा पॉलिसी खरीदी। बीमा कवरेज ₹2647217 है तथा वार्षिक प्रीमियम ₹17928 है। पॉलिसी 30/05/2025 से प्रभावी होगी ...
&
Financial / Insurance
\\

\hline

\end{tabular}

\end{table}

%%%%%%%%%%%%%%%%%%%%%%%%%%%%%%%%%%%%%%%%%%%%%%%%%%%%%%%%%%%%%%

\section{Annotation Strategy}
\label{sec:annotationstrategy}

\subsection{Existing Annotated Dataset}

Not applicable, since the dataset used in this project is self-created (see Section~\ref{subsec:selfcreated}) rather than sourced from an already-annotated third-party corpus.

%%%%%%%%%%%%%%%%%%%%%%%%%%%%%%%%%%%%%%%%%%%%%%%%%%%%%%%%%%%%%%

\subsection{Self-created Dataset}
\label{subsec:selfcreated}

\subsubsection*{Annotation Guidelines}

Each record was generated from one of seven fixed Hindi document templates (service agreement, court proceeding notice, loan document, insurance policy, tax notice, payment notice, confidentiality clause), where each template is intrinsically associated with exactly one \texttt{Document\_Type} and one \texttt{Intent} label. Populating a template with randomized entity values (party names, organization names, dates, amounts, and durations) therefore produces a record whose label is deterministically correct by construction.

\subsubsection*{Label Definitions}

\begin{itemize}
    \item \textbf{Legal -- Agreement:} Service agreements specifying parties, duration, and remuneration.
    \item \textbf{Legal -- Court Proceeding:} Notices describing case numbers, parties, and hearing dates.
    \item \textbf{Legal -- Confidentiality:} Clauses specifying non-disclosure obligations between parties.
    \item \textbf{Financial -- Finance:} General financial statements or notices.
    \item \textbf{Financial -- Loan:} Loan sanction documents specifying lender, borrower, amount, and interest rate.
    \item \textbf{Financial -- Payment:} Payment due or payment confirmation notices.
    \item \textbf{Financial -- Tax:} Tax-related notices or statements.
    \item \textbf{Financial -- Insurance:} Insurance policy documents specifying coverage and premium.
\end{itemize}

\subsubsection*{Annotation Process}

\begin{itemize}
    \item Number of annotators: Not applicable (labels assigned programmatically and deterministically at the time of template-based generation)
    \item Annotation tool used: Custom Python-based template-filling script
    \item Review process: Manual spot-checking of a random sample of generated records to confirm label correctness and text readability
\end{itemize}

\subsubsection*{Quality Checking}

\begin{itemize}
    \item Random sampling and manual review of generated records across all eight Intent categories
    \item Programmatic verification of zero missing values and zero duplicate \texttt{Text} entries
\end{itemize}

\subsubsection*{Inter-Annotator Agreement}

Not applicable, as labels are deterministically assigned by the template-generation process rather than by independent human annotators.

%%%%%%%%%%%%%%%%%%%%%%%%%%%%%%%%%%%%%%%%%%%%%%%%%%%%%%%%%%%%%%

\section{Dataset Licensing and Permissions}
\label{sec:datasetlicensing}

\begin{itemize}
    \item \textbf{License type:} Self-created for academic use; no external license applies.
    \item \textbf{Terms of use:} The dataset is used solely for the academic purposes of this Micro Project.
    \item \textbf{Citation requirements:} None, as the dataset is not derived from a third-party publication.
    \item \textbf{Redistribution policy:} The dataset may be shared for academic evaluation of this project.
    \item \textbf{Privacy considerations:} All party names, organization names, and case numbers appearing in the dataset are synthetically generated and do not correspond to real individuals, companies, or actual court cases.
    \item \textbf{Ethical considerations:} As the dataset is fully synthetic, no real personal or financial data is exposed, avoiding privacy and confidentiality concerns associated with real legal or financial documents.
\end{itemize}

%%%%%%%%%%%%%%%%%%%%%%%%%%%%%%%%%%%%%%%%%%%%%%%%%%%%%%%%%%%%%%

\section{Dataset Justification}

The selected dataset is suitable for the proposed project for the following reasons:

\begin{itemize}
    \item \textbf{Appropriate size:} 400 records provide a workable dataset for training and evaluating lightweight classifiers and fine-tuning transformer models within the scope of a micro project.
    \item \textbf{Relevant domain:} The dataset spans exactly the two domains -- Legal and Financial -- targeted by the proposed system.
    \item \textbf{Suitable language:} All text is in Hindi, matching the target language of the translation and classification stages.
    \item \textbf{Balanced classes:} A perfectly balanced 200/200 Document\_Type split, and reasonably balanced Intent classes (40--70 records each), reduce the risk of class-imbalance bias during model training.
    \item \textbf{High data quality:} No missing values or duplicate records were found during characteristic analysis.
    \item \textbf{Compatibility with the selected NLP task:} The dataset directly supports supervised training of the Document Classification module and provides realistic sentence structures (dates, monetary values, named entities) for evaluating the Key Information Extraction module.
\end{itemize}

\section{Final Dataset and Data Split}

The final intent dataset contains 400 Hindi records across nine intent labels. In the implemented prototype it is loaded as a complete reference set for query-intent signals; it is not used as evidence for answering questions about an uploaded document. For reproducible supervised experiments, the recommended stratified split is 70:15:15, with a fixed random seed and no template instance appearing in more than one partition.

\subsection{Training Dataset}
The proposed training partition contains 280 records (70\%). It is intended for fitting future statistical or transformer-based intent classifiers.

\subsection{Validation Dataset}
The proposed validation partition contains 60 records (15\%). It is reserved for threshold selection and model configuration, not final reporting.

\subsection{Testing Dataset}
The proposed testing partition contains 60 records (15\%). It is held out for final intent-classification evaluation. Separately, the current Review 3 system evaluation uses eight labelled bilingual PDF documents to assess document-domain classification and language detection.

%%%%%%%%%%%%%%%%%%%%%%%%%%%%%%%%%%%%%%%%%%%%%%%%%%%%%%%%%%%%%%

\section*{Chapter Summary}

This chapter presented the dataset selected for the proposed \textit{DocSaarthi} NLP project: a self-created, 400-record, template-generated Hindi legal and financial text corpus with balanced Document\_Type and Intent labels. The dataset source, description, characteristics, annotation strategy, licensing information, and suitability for the selected NLP task were discussed. A well-documented dataset provides the foundation for effective preprocessing and model development. The next chapter describes the preprocessing techniques and exploratory data analysis performed on the collected dataset.
